\section{Introduction}
\label{sec:intro}

In the legal domain, statutes are a fundamental source of knowledge that guides principles of jurisdiction \citep{ilpcsr}. In practice, when faced with a legal situation, a practitioner identifies applicable statutes based on experience and knowledge. This process is time-consuming, and the problem worsens with the growing number of legal cases in populous countries. Hence, there is an imminent need to automate statute retrieval: given a description of facts, identify the relevant legal provisions.

Statute retrieval requires bridging two fundamentally different registers. Queries describe concrete events in narrative language; statutes define abstract legal norms using formal terminology. A layperson asking about a landlord who refuses to return a deposit must be matched to provisions on \textit{obligation} (\textit{perikatan}) and \textit{debtor default} (\textit{wanprestasi}). This vocabulary gap is the core challenge, and it widens as queries become less formal.

Current statute retrieval benchmarks do not reveal whether methods are robust to this challenge. IL-PCSR \citep{ilpcsr} uses full case judgments and LLM-generated summaries as queries. COLIEE \citep{coliee2024} provides expert-crafted bar exam questions. BSARD \citep{bsard} and STARD \citep{stard} use shorter queries, but each evaluates at a single formality level. While STARD demonstrated that models trained on formal queries perform poorly on real layperson queries, their comparison is cross-dataset rather than controlled: different queries about different cases make it impossible to isolate the effect of query formality. Additionally, no existing work evaluates whether a single retrieval method generalizes across multiple languages and legal systems. We address these gaps with the following contributions:

\begin{itemize}
\item We construct KUHPerdata, a statute retrieval dataset for Indonesian civil law comprising court cases mapped to 2,127 articles from the \textit{Kitab Undang-Undang Hukum Perdata} (Indonesian Civil Code). The dataset includes two query variants generated from the same source court decisions: (i) narrative case summaries and (ii) plain-language layperson questions, sharing identical relevance judgments. To the best of our knowledge, this is the first statute retrieval resource for Indonesian, a language spoken by 270 million people.

\item We conduct error analysis across four multilingual datasets (Indonesian, French, Chinese, English) and identify two consistent patterns: (a) vocabulary overlap is the primary predictor of lexical retrieval failure across all four languages, and (b) co-relevant statutes cluster structurally (median 18 articles apart, 85\% in the same legal domain) while hard negatives are distant (median 297 articles). We also observe that the optimal balance between lexical and semantic retrieval signals varies by dataset and correlates with the level of lexical overlap.

\item Based on these findings, we extend the paragraph-level graph neural network of \citet{ilpcsr} with two modifications: statute proximity edges that connect structurally nearby articles in the retrieval graph, and post-training alpha grid search that adapts the lexical-semantic blend per dataset. We call this extension Prox-GNN. On four datasets, Prox-GNN achieves MRR@10 of 0.486--0.527, improving 24--94\% over the strongest baseline. Ablation confirms that each modification addresses the specific failure pattern that motivated it.
\end{itemize}
